\centerline{\bf \Huge Зацикливания}

\begin{flushright}
        \scriptsize{Тот, кто зацикливается на чём-то, теряет всё остальное\\
        Монолог фармацевта 2 сезон 9 серия}
\end{flushright}

\textbf{Принцип зацикливания.} Если система может находиться лишь в конечном числе состояний, и каждое следующее состояние однозначно определяется по предыдущему, то система с некоторого момента зациклится.

\textbf{Принцип зацикливания назад.} Если в условиях принципа зацикливания каждое предыдущее состояние определяется по следующему то система зацикливается без предпериода.



\problem{1}
В некоторой стране стали очень популярны путешествия. При этом часть авиарейсов прекратило свою работу так, что оказалось, что теперь из каждого города ведёт ровно один односторонний маршрут в другой город. При этом в каждый город ведёт не более одного маршрута. Докажите, что турист, выехавший из столицы, когда-нибудь в неё вернётся.

\problem{2}
В стране 16 городов, некоторые из них соединены дорогой с односторонним движением, в каждый город ровно одна дорога входит и из каждого города ровно одна дорога выходит. В каждом городе находится по автомобилисту. Каждый день каждый автомобилист проезжает одну дорогу. Найдите наименьшее натуральное $N$ такое, что ровно через $N$ дней все автомобилисты будут в тех городах, из которых начинали движение, вне зависимости от того, как именно проложены дороги.

\problem{3}
Метеорологическая служба Цветочного города следит за погодой уже сто лет. Они подразделяют погоду на дождливую или солнечную. Метеорологи доказали, что погода на следующий день однозначно определяется (по какому-то загадочному принципу) погодой в предыдущие семь дней. Последняя неделя в Цветочном городе была полностью солнечная. Докажите, что в будущем снова встретится полностью солнечная неделя.


\problem{4}
Напишем последовательность $2,0,2,3,7,2,4,6, \ldots$, в которой каждый новый член равен последней цифре суммы четырёх предыдущих. Докажите, что рано или поздно в последовательности встретится кусок

(a) $2,0,2,3$;

(б) $4,8,1,9$.

\newpage

\hspace{3mm}

(в) Докажите, что количество членов последовательности между повторяющимися кусками $2,0,2,3$ (считая первый кусок $2,0,2,3$ и не считая второй) делится на 5.

(г) Докажите, что кусок $2,0,2,3$ встретится не позже чем через $5^5$ членов в последовательности.

\problem{5}
Кубик Рубика вывели из исходного состояния некоторой последовательностью поворотов граней. Докажите, что если повторять эту последовательность поворотов достаточно долго, то кубик в конце концов вернется в исходное состояние.